% Input fragment: the containing document supplies packages and notation macros.
\section{Design zeros, information blocks, and what an inverse means}
\label{cov-structure}

\subsection{Dimensions and the statistical object being inverted}
\label{cov-dimensions}

Throughout this fragment, there are $M$ studies, $N$ voxels, $P$ spatial
basis coefficients per map, $G$ groups, $R$ spatially varying (SV)
covariates, and $Q$ global covariates.  The basis matrix is
$B\in\R^{N\times P}$, with row $j$ written $B_j^\top$:
$B_j\in\R^P$ is a column vector.  Likewise $Z_i\in\R^R$ and
$z_i\in\R^Q$ are column vectors of study covariates.  The coefficient
vector $\theta$ contains regression coefficients only unless dispersion
is explicitly added in Section~\ref{cov-dispersion}.

For SV maps, write
$\beta=(\beta_1^\top,\ldots,\beta_R^\top)^\top\in\R^{RP}$.
For group maps, use
$\beta=(\xi_1^\top,\ldots,\xi_G^\top)^\top\in\R^{GP}$.
Thus the symbol $\beta$ always denotes stacked spatial coefficients,
but its dimension depends on the specification.  With global coefficients,
$\theta=(\beta^\top,\gamma^\top)^\top$ and $\gamma\in\R^Q$.
Let $d=\dim(\beta)$ and $n_\theta=d+Q$; set $Q=0$ when no global
coefficients are fitted.  Typical affine predictors, with a known offset
$o_{ij}$, are
\begin{align}
 \eta_{ij}
 &=o_{ij}+(Z_i^\top\otimes B_j^\top)\beta+z_i^\top\gamma
 &&\text{(SV specification)},\\
 \eta_{ij}
 &=o_{ij}+B_j^\top\xi_{g(i)}+z_i^\top\gamma
 &&\text{(group specification)}.
\end{align}
An offset is not an additional estimated coefficient.

Our convention is
\begin{equation}
 H(\theta)=-\nabla_\theta^2\ell(\theta),\qquad
 I(\theta)=\E_\theta\{H(\theta)\}.
 \label{cov-information-convention}
\end{equation}
Here $H$ is \emph{observed} information.  Under a correctly specified
regular likelihood, interchange of differentiation and integration also
gives $I=\Var_\theta\{\nabla_\theta\ell\}$.  That information equality
is a statistical assertion with hypotheses, not a consequence of matrix
symmetry.  Dispersion $\alpha$ is fixed for now.

Identities below about $H^{-1}$ are exact finite-dimensional algebra when
the stated inverses exist.  Interpreting $H(\widehat\theta)^{-1}$ or
$I(\widehat\theta)^{-1}$ as a sampling covariance is generally a
first-order likelihood approximation, requiring identification, a suitable
independent-unit asymptotic regime, and an interior regular fit.
An indefinite observed information matrix may be invertible without its
inverse being a valid covariance.  These two questions must remain separate.

\subsection{From coefficient support to a graph of SV blocks}
\label{cov-support-graph}

For an independent-cell log likelihood, put
$w_{ij}=-\partial^2\ell_{ij}/\partial\eta_{ij}^2$.
The affine-design chain rule, rather than any particular count family,
gives the $P\times P$ spatial block
\begin{equation}
 H_{rr'}=\sum_{i=1}^M\sum_{j=1}^N
             w_{ij}Z_{ir}Z_{ir'}B_jB_j^\top .
 \label{cov-sv-block}
\end{equation}
If $Z_{ir}Z_{ir'}=0$ for every study $i$, every summand vanishes,
irrespective of the numerical weights.  Consequently $H_{rr'}=0$.
Conversely, a zero block at one parameter value might arise from
cancellation of nonzero terms.  Such an accidental zero is not a
design-guaranteed decomposition that can safely be reused at other fits.

Construct a graph with one vertex for each SV covariate column.  Connect
$r$ and $r'$ whenever at least one study has both $Z_{ir}\ne0$ and
$Z_{ir'}\ne0$.  If two vertices lie in different connected components,
there is no such study, and their entire information cross block vanishes.
If component $a$ has $r_a$ columns, it contributes $d_a=r_aP$
coefficients.  After permuting coefficients by components,
\begin{equation}
 D=\blkdiag(D_1,\ldots,D_{K_{\mathrm{comp}}}),\qquad
 \sum_{a=1}^{K_{\mathrm{comp}}}d_a=d.
 \label{cov-component-blocks}
\end{equation}
For a group-only spatial design, these are precisely the $G$ group
blocks, so we write $D=\blkdiag(D_g:g=1,\ldots,G)$.
For a general SV design, component labels $a$ replace group labels $g$;
a component need not be a scientific group.

A path is enough to prevent independent inversion.  For example, suppose
column 1 co-occurs with column 2, and column 2 with column 3, while
columns 1 and 3 never co-occur.  The information can have the form
\[
 \begin{pmatrix} a&b&0\\b&c&e\\0&e&f\end{pmatrix}.
\]
When it is nonsingular, its $(1,3)$ inverse entry is
$be/\det(H)$, which is generally nonzero.  A missing direct edge does
not imply a zero inverse entry.  Disjoint cell-means indicators separate;
a common SV intercept or another overlapping column can connect them.
An all-zero design column instead produces an unidentified direction.

For a likelihood factorized by study but not by voxel, a study contribution
may couple its voxels.  Nevertheless, it cannot differentiate with respect
to two spatial columns that are not both present in that study, so the
same support argument applies when that factorization and predictor support
hold.  If an aggregate likelihood couples studies or groups, its actual
factor supports must be inspected instead.  Similarly, a penalty that
couples components adds edges; design zeros for the likelihood alone
are not then zeros for the penalized objective.

\subsection{Verification of independent block inversion}
\label{cov-block-verification}

First suppose there is no global border.  Let $\Pi$ be the permutation
matrix putting the components in order, and suppose each $D_a$ is
nonsingular.  Define
\[
 W=\Pi^\top\blkdiag(D_1^{-1},\ldots,D_{K_{\mathrm{comp}}}^{-1})\Pi,
 \qquad H=\Pi^\top D\Pi .
\]
Since $\Pi\Pi^\top=I_d$, multiplication gives
\[
 HW
 =\Pi^\top
   \blkdiag(D_1D_1^{-1},\ldots,D_{K_{\mathrm{comp}}}
                                  D_{K_{\mathrm{comp}}}^{-1})\Pi
 =\Pi^\top I_d\Pi=I_d .
\]
Multiplying in the reverse order gives $WH=I_d$ as well.  Thus $W=H^{-1}$.
This is an exact identity, not an approximation obtained by ignoring small
cross blocks.  It also shows why parameter ordering must accompany a
collection of blocks: block coordinates and original coefficient
coordinates need not coincide.

With shared global coefficients, the full information instead has the form
\begin{equation}
 H=\begin{pmatrix}D&F\\F^\top&\mathcal E\end{pmatrix},
 \quad F\in\R^{d\times Q},\quad\mathcal E\in\R^{Q\times Q}.
 \label{cov-bordered-information}
\end{equation}
Here $\mathcal E$ denotes a matrix, not an expectation or an intensity
total.  Although $D$ has disconnected spatial blocks, paths through the
global coefficients connect the full system.

\section{Schur elimination and the covariance induced by shared coefficients}
\label{cov-schur}

\subsection{Deriving a solve one equation at a time}
\label{cov-schur-solve}

Assume $D$ is nonsingular.  For a right-hand side
$b=(b_D^\top,b_\gamma^\top)^\top$, write the desired solution as
$x=(x_D^\top,x_\gamma^\top)^\top$.  The two block equations in $Hx=b$ are
\begin{align}
 Dx_D+Fx_\gamma&=b_D, \label{cov-first-block-equation}\\
 F^\top x_D+\mathcal E x_\gamma&=b_\gamma .
\end{align}
Solving the first equation for the spatial part gives
$x_D=D^{-1}b_D-D^{-1}Fx_\gamma$.  Substitute this expression into
the second equation:
\begin{align}
 F^\top D^{-1}b_D-F^\top D^{-1}Fx_\gamma
                      +\mathcal E x_\gamma&=b_\gamma,\\
 (\mathcal E-F^\top D^{-1}F)x_\gamma
                      &=b_\gamma-F^\top D^{-1}b_D .
\end{align}
Define, once for all right-hand sides,
\begin{equation}
 T=D^{-1}F,\qquad S=\mathcal E-F^\top T .
 \label{cov-schur-definitions}
\end{equation}
If $S$ is nonsingular, the complete solve is therefore
\begin{align}
 y&=D^{-1}b_D,\\
 x_\gamma&=S^{-1}(b_\gamma-F^\top y),\\
 x_D&=y-Tx_\gamma .
 \label{cov-schur-algorithm}
\end{align}
The notation $D^{-1}b_D$ means a solve with the factorized blocks of $D$;
it is not an instruction to construct a dense inverse.  These equations
also hold for several right-hand sides stored as columns.

\subsection{Full inverse and explicit multiplication check}
\label{cov-full-inverse}

Expanding the solution in its two right-hand-side components gives
\[
 x_D=(D^{-1}+TS^{-1}T^\top)b_D-TS^{-1}b_\gamma,\qquad
 x_\gamma=-S^{-1}T^\top b_D+S^{-1}b_\gamma .
\]
We used symmetry of $D$ to obtain $F^\top D^{-1}=T^\top$.
Hence
\begin{equation}
 H^{-1}=
 \begin{pmatrix}
 D^{-1}+TS^{-1}T^\top&-TS^{-1}\\
 -S^{-1}T^\top&S^{-1}
 \end{pmatrix}.
 \label{cov-inverse-formula}
\end{equation}
For a direct audit, multiply $H$ by the proposed inverse.  Its four
blocks simplify separately, using $DT=F$:
\begin{align}
 (HH^{-1})_{DD}
 &=I_d+DTS^{-1}T^\top-FS^{-1}T^\top=I_d,\\
 (HH^{-1})_{D\gamma}
 &=-DTS^{-1}+FS^{-1}=0,\\
 (HH^{-1})_{\gamma D}
 &=T^\top+(F^\top T-\mathcal E)S^{-1}T^\top
   =T^\top-SS^{-1}T^\top=0,\\
 (HH^{-1})_{\gamma\gamma}
 &=(\mathcal E-F^\top T)S^{-1}=I_Q .
\end{align}
The transpose verifies the reverse product because both matrices are
symmetric.  When $F=0$, the formula reduces to separate inverses of
$D$ and $\mathcal E$.  When $Q=0$, there is no Schur calculation.

Invertibility of $H$ alone does not guarantee invertibility of $D$.
For example,
$\left(\begin{smallmatrix}0&1\\1&0\end{smallmatrix}\right)$
is nonsingular but its leading scalar block is zero.  The particular
elimination above is unavailable in that case; a different pivoting
scheme or full symmetric indefinite factorization is needed.

\subsection{Positive definiteness by completing the square}
\label{cov-schur-positive}

For arbitrary $u\in\R^d$ and $v\in\R^Q$, expand
\begin{align}
 \begin{pmatrix}u\\v\end{pmatrix}^{\!\top}
 H\begin{pmatrix}u\\v\end{pmatrix}
 &=u^\top Du+2u^\top Fv+v^\top\mathcal E v\\
 &=(u+Tv)^\top D(u+Tv)+v^\top Sv .
 \label{cov-completed-square}
\end{align}
Indeed, the cross term is $2u^\top DTv=2u^\top Fv$ and the
remaining $v$ term is
$v^\top(T^\top DT+S)v=v^\top\mathcal E v$.
Thus, if $D\succ0$, then $H\succ0$ if and only if $S\succ0$.
Necessity follows by choosing $u=-Tv$; sufficiency follows because
the sum is positive for any nonzero $(u,v)$.  With $D\succ0$,
the analogous statement $H\succeq0$ if and only if $S\succeq0$
holds, but a singular $S$ prevents an ordinary inverse.

Equivalently,
\begin{equation}
 H=
 \begin{pmatrix}I_d&0\\T^\top&I_Q\end{pmatrix}
 \begin{pmatrix}D&0\\0&S\end{pmatrix}
 \begin{pmatrix}I_d&T\\0&I_Q\end{pmatrix}.
 \label{cov-congruence}
\end{equation}
This is a congruence, not an orthogonal diagonalization.  It determines
inertia and gives $\det(H)=\det(D)\det(S)$, but it does not say that
the eigenvalues of $D$ and $S$ are the eigenvalues of $H$.

\subsection{Joint versus fixed-parameter covariance and profile curvature}
\label{cov-joint-conditional}

Suppose $H\succ0$ and use $V=H^{-1}$ as the local model-based
covariance approximation.  If $\gamma$ were known rather than estimated,
the spatial curvature would be $D$ and its covariance approximation
would be $D^{-1}$.  Joint estimation instead gives
\begin{equation}
 V_{\beta\beta}=D^{-1}+TS^{-1}T^\top,\qquad
 V_{\beta\gamma}=-TS^{-1},\qquad V_{\gamma\gamma}=S^{-1}.
 \label{cov-joint-blocks}
\end{equation}
For any $a\in\R^d$,
$a^\top TS^{-1}T^\top a=(T^\top a)^\top S^{-1}(T^\top a)\geq0$.
Thus the extra joint-estimation term is positive semidefinite (PSD).
This statement is about inverse curvature; it is not a general
ordering theorem for arbitrary robust covariance estimators.

In a Gaussian approximation with precision $H$, completing the square
also shows that the conditional covariance of $\beta$ given $\gamma$
is $D^{-1}$.  This is the precise local meaning of ``conditional'' here;
it does not assert an exact conditional sampling law for an arbitrary
finite-sample estimator.  Similarly, $\mathcal E^{-1}$ is the local
covariance for $\gamma$ if $\beta$ were known.  Joint inference uses
$S^{-1}$ instead.  Since
$\mathcal E-S=F^\top D^{-1}F\succeq0$, both matrices being SPD implies
$S^{-1}\succeq\mathcal E^{-1}$.

The Schur complement also has an optimization interpretation.  Let
$L=-\ell$, and let $\beta_*(\gamma)$ be a differentiable local minimizer
in $\beta$, satisfying $L_\beta(\beta_*(\gamma),\gamma)=0$.
Implicit differentiation gives
\[
 D\,\frac{\partial\beta_*}{\partial\gamma^\top}+F=0,
 \qquad \frac{\partial\beta_*}{\partial\gamma^\top}=-D^{-1}F=-T .
\]
For the profile objective $L_{\mathrm{prof}}(\gamma)=
L(\beta_*(\gamma),\gamma)$, the term involving $L_\beta$ vanishes:
\[
 \nabla_\gamma L_{\mathrm{prof}}=L_\gamma,\qquad
 \nabla_\gamma^2L_{\mathrm{prof}}
 =\mathcal E+F^\top\frac{\partial\beta_*}{\partial\gamma^\top}
 =\mathcal E-F^\top D^{-1}F=S .
\]
These derivatives are exact along a smooth stationary branch.  Treating
this branch as a local profile minimum additionally requires appropriate
curvature; global uniqueness does not follow from the calculation.

\subsection{Cross-group covariance and a hand-computable example}
\label{cov-example}

For group blocks let $T_g=D_g^{-1}F_g$, so $T$ stacks the $T_g$.
Equation~\eqref{cov-joint-blocks} gives
\begin{equation}
 V_{\xi_g\xi_h}=\mathbf1\{g=h\}D_g^{-1}+T_gS^{-1}T_h^\top .
 \label{cov-cross-group}
\end{equation}
The entire spatial correction $TS^{-1}T^\top$ has rank at most $Q$.
Each cross-group block has rank at most $Q$ as well, though its entries
can be positive or negative; an off-diagonal rectangular block is not
itself a PSD matrix.  A small number of global coefficients therefore
induces low-rank, not necessarily negligible, dependence among many maps.

Take $G=2$, $P=Q=1$, and choose
\[
 D=\begin{pmatrix}2&0\\0&3\end{pmatrix},\qquad
 F=\begin{pmatrix}1\\1\end{pmatrix},\qquad
 \mathcal E=2.
\]
The arithmetic in the elimination is
\[
 T=\begin{pmatrix}1/2\\1/3\end{pmatrix},\quad
 F^\top T=\frac12+\frac13=\frac56,\quad
 S=2-\frac56=\frac76,\quad S^{-1}=\frac67.
\]
Both $D$ and $S$ are positive definite, so the bordered matrix is SPD.
The spatial correction is
\[
 TS^{-1}T^\top=
 \frac67\begin{pmatrix}1/4&1/6\\1/6&1/9\end{pmatrix}
 =\begin{pmatrix}3/14&1/7\\1/7&2/21\end{pmatrix}.
\]
Adding $D^{-1}$ and the remaining blocks yields
\begin{equation}
 H=\begin{pmatrix}2&0&1\\0&3&1\\1&1&2\end{pmatrix},
 \qquad
 H^{-1}=\frac17
 \begin{pmatrix}5&1&-3\\1&3&-2\\-3&-2&6\end{pmatrix}.
 \label{cov-numerical-inverse}
\end{equation}
For example, the first row of the product is
$\{2(5,1,-3)+(-3,-2,6)\}/7=(1,0,0)$; the other rows
are checked in the same way.  Spatial estimates have covariance $1/7$
despite a zero spatial information cross entry.

For the group difference $c^\top\theta=\xi_1-\xi_2$,
\[
 c^\top H^{-1}c=\frac57+\frac37-2\frac17=\frac67.
\]
Using the correct marginal variances but discarding their covariance
would give $8/7$, not $6/7$.  Treating $\gamma$ as known would instead
give $1/2+1/3=5/6$.  The joint correction for this difference is only
$(1/2-1/3)^2(6/7)=1/42$, and indeed $5/6+1/42=6/7$.
This demonstrates both distinctions: joint versus known global effects,
and joint marginals versus incorrectly independent marginals.  The
direction of error from dropping cross covariances need not always be
the same as in this example.

\section{Selected covariance calculations without a full inverse}
\label{cov-selected}

\subsection{Project by solving for the contrasts}
\label{cov-projection}

Let $C\in\R^{k\times n_\theta}$ contain $k$ fixed linear contrasts.
Solving
\begin{equation}
 HX=C^\top,\qquad V_C=CX
 \label{cov-selected-solve}
\end{equation}
gives $X=H^{-1}C^\top$ and hence $V_C=CH^{-1}C^\top$.
Thus only $k$ right-hand sides are needed, rather than $n_\theta$ columns
of an inverse.  Symmetry gives
$V_C^\top=(CH^{-1}C^\top)^\top=V_C$ exactly; numerical
symmetrization can remove roundoff asymmetry but cannot repair a wrong
model or an indefinite result.

There is a useful reduced formula.  Partition $C=[\,C_D\ C_\gamma\,]$
conformably with the border.  Multiplication of the four inverse blocks
gives
\begin{align}
 V_C={}&C_DD^{-1}C_D^\top+C_DTS^{-1}T^\top C_D^\top\\
 &-C_DTS^{-1}C_\gamma^\top
       -C_\gamma S^{-1}T^\top C_D^\top
       +C_\gamma S^{-1}C_\gamma^\top\\
 ={}&C_DD^{-1}C_D^\top+
       (C_\gamma-C_DT)S^{-1}(C_\gamma-C_DT)^\top .
 \label{cov-reduced-contrast}
\end{align}
For a spatial-only contrast, its uncertainty from the global fit vanishes
exactly when $C_DT=0$ (assuming $S\succ0$).  For a contrast involving
both kinds of coefficients, cancellation instead concerns
$C_\gamma-C_DT$.  These statements follow from the positive definite
quadratic form and are not an assumption of group independence.

Computationally, solve $DY=C_D^\top$, form
$A_C=C_\gamma-C_DT$, and solve $SW=A_C^\top$.  Then compute
$C_DY+A_CW$.  Neither $D^{-1}$ nor $S^{-1}$ must be stored as a
matrix.  The same projection principle works with an expected information
matrix or a specified penalized curvature, but that substitution changes
which covariance interpretation is being used.

\subsection{Diagonal entries and contrasts in batches}
\label{cov-diagonal}

Let $e_a$ denote a coordinate vector and $t_a=T^\top e_a\in\R^Q$.
For a spatial coordinate in block $g$,
\begin{equation}
 (H^{-1})_{aa}=(D_g^{-1})_{aa}+t_a^\top S^{-1}t_a .
 \label{cov-diagonal-entry}
\end{equation}
The first term means the corresponding local coordinate in $D_g$.
To compute the second terms together, solve $SW=T^\top$;
the $a$th correction is the inner product of row $a$ of $T$
with column $a$ of $W$.  Global diagonal entries are the diagonal
of $S^{-1}$, obtained from solves against its coordinate vectors.

Even obtaining every diagonal entry of a dense block inverse may require
many triangular solves.  The benefit is that the full $d\times d$
spatial covariance is never formed.  If just a few coordinates are
required, solve only for their coordinate vectors.  If many voxel
contrasts are needed, process a manageable batch of contrast rows and
retain the resulting variances or small covariance blocks.  An
$N\times N$ covariance of all voxels should be formed only if it is
actually the requested output.

\section{Independent-unit sandwich covariance and its normalization}
\label{cov-sandwich}

\subsection{Taylor expansion of an unnormalized score equation}
\label{cov-sandwich-taylor}

Let $U_i(\theta)\in\R^{n_\theta}$ be the contribution of independent
study $i$ to an estimating equation, and let
$U(\theta)=\sum_{i=1}^M U_i(\theta)$.  At the target $\theta_0$ assume
$\E\{U(\theta_0)\}=0$; common likelihood and conditional-mean setups
also have $\E\{U_i(\theta_0)\}=0$ separately.  Define
\[
 A(\theta)=-\frac{\partial U(\theta)}{\partial\theta^\top}
          =-\sum_i\frac{\partial U_i(\theta)}{\partial\theta^\top}.
\]
Expanding $U(\widehat\theta)=0$ about $\theta_0$ gives
\begin{align}
 0&=U(\theta_0)-A(\theta_0)(\widehat\theta-\theta_0)
                      +\text{higher-order remainder},\\
 \widehat\theta-\theta_0
 &\approx A(\theta_0)^{-1}\sum_iU_i(\theta_0).
 \label{cov-taylor-linearization}
\end{align}
This linearization requires a negligible remainder in the chosen
asymptotic regime and a nonsingular limiting sensitivity.  Independence
between studies gives
\[
 \Var\Bigl\{\sum_iU_i(\theta_0)\Bigr\}
 =\sum_i\Var\{U_i(\theta_0)\}=J_0 .
 \]
With appropriate laws of large numbers, a customary empirical replacement
is $\widehat J=\sum_iU_i(\widehat\theta)U_i(\widehat\theta)^\top$,
leading to
\begin{equation}
 \widehat V_{\mathrm{robust}}
 =\widehat A^{-1}\widehat J\,\widehat A^{-\top}.
 \label{cov-sandwich-estimator}
\end{equation}
Here $A^{-\top}=(A^{-1})^\top$.  Likelihood scores give $A=H$;
general estimating equations need not give a symmetric $A$.
Consistency of the uncentered meat requires the relevant moment
conditions.  In particular, systematic nonzero individual score means
under a misspecified fixed design cannot simply be identified with
within-study sampling variances.

There is \emph{no additional factor $1/M$} in
Equation~\eqref{cov-sandwich-estimator}: its bread and meat are both
unnormalized sums.  To check, write
$\bar A=\widehat A/M$ and $\bar J=\widehat J/M$.  Then
\[
 \frac1M\bar A^{-1}\bar J\bar A^{-\top}
 =\frac1M(M\widehat A^{-1})(\widehat J/M)
                     (M\widehat A^{-\top})
 =\widehat A^{-1}\widehat J\widehat A^{-\top}.
\]
Mixing an averaged bread with an unscaled summed meat introduces
spurious powers of $M$.

\subsection{Why the clustering unit matters}
\label{cov-study-clustering}

Suppose a working likelihood is cell additive and
$U_i=\sum_{j=1}^N u_{ij}$.  Expanding its outer product shows exactly
what study clustering retains:
\begin{align}
 U_iU_i^\top
 &=\Bigl(\sum_j u_{ij}\Bigr)\Bigl(\sum_lu_{il}\Bigr)^\top\\
 &=\sum_j u_{ij}u_{ij}^\top
       +\sum_{j\ne l}u_{ij}u_{il}^\top .
 \label{cov-cluster-cross-terms}
\end{align}
A cellwise meat retains only the first sum.  A study-clustered meat
retains both, allowing dependence among voxels within a study.  It still
requires independent studies, or a justified asymptotic replacement for
that assumption.  If experiments share participants or another sampling
unit, the independence unit may be larger than an experiment and the
scores must be grouped accordingly.

More voxels within a fixed small number of studies do not automatically
provide the independent replication required by cluster asymptotics.
Finite-sample adjustments also do not create that replication.  For
example, multiplying $\widehat J$ by $M/(M-n_\theta)$ is a particular
HC1 convention, defined only for $M>n_\theta$; it is not a universal
small-cluster validity theorem.

\subsection{Family-specific scores cannot be interchanged}
\label{cov-family-scores}

Write $x_{ij}=\nabla_\theta\eta_{ij}$ as a column vector.  For the
Poisson working likelihood its study score is
\[
 U_i^{\mathrm P}=\sum_jx_{ij}(y_{ij}-\mu_{ij}).
\]
For independent-cell NB2 with fixed dispersion $\alpha$, the corresponding
score is instead
\[
 U_i^{\mathrm{NB}}=\sum_jx_{ij}
                  \frac{y_{ij}-\mu_{ij}}{1+\alpha\mu_{ij}} .
\]
These formulas state the score inputs; the count-likelihood derivative
derivations belong to the likelihood chapters.  The NB denominators
change both the estimating equation and its derivative.  Combining
Poisson residual scores with NB likelihood information is not the
matching likelihood sandwich.  Under a correctly specified conditional
mean, the expected derivative of this independent-NB working score
has the usual NB Fisher weights, even if the working variance is wrong;
using those expected weights as bread then needs its own consistency
argument.  Observed NB information and expected Fisher information do
not in general agree at a realized dataset.

For the SV predictor the Poisson spatial and global score pieces are
explicitly
\[
 U_{i,\beta}^{\mathrm P}
    =Z_i\otimes\left\{\sum_jB_j(y_{ij}-\mu_{ij})\right\},
 \qquad
 U_{i,\gamma}^{\mathrm P}
    =z_i\sum_j(y_{ij}-\mu_{ij}).
\]
Thus scores can be assembled by projecting residuals onto the basis,
without constructing every full design row.

A study-level clustered-NB likelihood obtained by integrating a shared
study effect has score $U_i=\nabla_\theta\ell_i$, including that
integration-induced dependence; it is not an independent-cell NB score.
An aggregate-NB likelihood for group-by-voxel totals is more problematic
for study clustering: its objective need not be a sum of original-study
log likelihoods at all.  A matched shape parameter can depend on several
studies' means simultaneously.  Arbitrarily allocating an aggregate
score back to studies does not supply independent score contributions,
even if the allocated vectors sum to the total score.  A valid sandwich
must follow an actual independent-unit estimating representation or a
justified joint sampling model for the aggregates.  Alternatively, a
bootstrap can refit the aggregate procedure to appropriately resampled
original independent units when those data and assumptions are available.

In \texttt{cbmr\_improved/covariance.py}, the implemented likelihood-score
sandwich explicitly requires Poisson and rejects the NB families.
Its study-score path uses basis projections of counts and fitted means.
The matrix machinery in \texttt{structured.py} can accept other supplied
scores, but cannot certify that the caller's scores represent the fitted
likelihood or the correct independence units.

\subsection{Projected Gram computation and rank limits}
\label{cov-sandwich-gram}

Let $\mathsf U\in\R^{M\times n_\theta}$ have row $i$ equal to
$U_i^\top$, so $\widehat J=\mathsf U^\top\mathsf U$.
For symmetric bread $H$ and a contrast matrix $C$, solve
$HX=C^\top$ and put $\mathsf V=\mathsf U X$.  Then
\begin{align}
 \mathsf V^\top\mathsf V
 &=X^\top\mathsf U^\top\mathsf U X\\
 &=CH^{-1}\widehat JH^{-1}C^\top
 =C\widehat V_{\mathrm{robust}}C^\top .
 \label{cov-projected-gram}
\end{align}
Equivalently, accumulate outer products
$v_iv_i^\top$ with $v_i=CH^{-1}U_i$.
For a nonsymmetric bread, the correct projected solve is
$A^\top X=C^\top$, not $AX=C^\top$.
This distinction disappears for the symmetric observed information.

The Gram representation proves PSD of the sandwich even if an
invertible bread is indefinite.  It does not make such a bread a valid
likelihood covariance or prove consistency of the estimator.  It also
shows
\[
 \operatorname{rank}(\widehat V_{\mathrm{robust}})
 \leq\min(M,n_\theta),\qquad
 \operatorname{rank}(C\widehat V_{\mathrm{robust}}C^\top)
 \leq\min(M,n_\theta,k).
\]
At an exactly solved unpenalized score equation,
$\mathsf U^\top\mathbf1_M=\sum_iU_i=0$, which tightens the first
bound to $\min(M-1,n_\theta)$ and similarly for the projected matrix.
Explicitly centering score rows also imposes this dependence.
At a penalized optimum the unpenalized score sum is generally nonzero,
so the $M-1$ argument does not automatically apply.  Multiplying a
Gram matrix by a positive correction factor cannot increase its rank.

\section{Quadratic roughness penalties and the meaning of penalized uncertainty}
\label{cov-penalties}

\subsection{Deriving the gradient and Hessian from the quadratic form}
\label{cov-penalty-derivatives}

Let $J_B=J_B^\top\succeq0$ be a fixed $P\times P$ roughness matrix
and let $\lambda_g\geq0$ be fixed smoothing parameters.  For group maps
use the explicit convention
\begin{equation}
 \mathcal P(\theta)=\frac12\sum_{g=1}^G
                         \lambda_g\xi_g^\top J_B\xi_g
                   =\frac12\theta^\top K\theta,\quad
 K=\blkdiag(\lambda_1J_B,\ldots,\lambda_GJ_B,0_{Q\times Q}).
 \label{cov-penalty-definition}
\end{equation}
For SV maps, replace the group blocks by
$\lambda_rJ_B$, $r=1,\ldots,R$.  The construction of $J_B$ depends
on the specified roughness operator and basis; it is not determined by
the word ``smoothing'' alone.

For one coefficient vector $\xi$, expand
$\xi^\top J_B\xi=\sum_{a,b}\xi_a(J_B)_{ab}\xi_b$.
Differentiation with respect to $\xi_c$ gives
\begin{align}
 \frac{\partial}{\partial\xi_c}
       \left(\frac{\lambda}{2}\xi^\top J_B\xi\right)
 &=\frac{\lambda}{2}
       \left\{\sum_b(J_B)_{cb}\xi_b+\sum_a\xi_a(J_B)_{ac}\right\}\\
 &=\lambda(J_B\xi)_c .
\end{align}
Differentiating again gives
$\partial^2\mathcal P/\partial\xi_c\partial\xi_e
=\lambda(J_B)_{ce}$.  Equivalently, the differential is
$d\mathcal P=(K\theta)^\top d\theta$ and its second differential
is $d\theta^\top K\,d\theta$.  A nonsymmetric supplied matrix would
contribute its symmetric part; the convention here assumes symmetry.

For the minimized objective $L_{\mathrm{pen}}=-\ell+\mathcal P$,
\begin{equation}
 \nabla L_{\mathrm{pen}}=-U+K\theta,\qquad
 H_{\mathrm{pen}}=H+K,\qquad U(\widehat\theta)=K\widehat\theta .
 \label{cov-penalized-equation}
\end{equation}
Omitting the factor $1/2$ in the penalty doubles both derivatives at the
same numerical $\lambda_g$.  Fixed separate-map penalties preserve
the likelihood's spatial component zeros and global border.  A
between-map penalty need not preserve them.

\subsection{Frequentist linearization, bias, and posterior curvature}
\label{cov-penalty-uncertainty}

Let $\theta_0$ be the unpenalized population target.  Expanding the
penalized score equation about that target gives
\begin{align}
 0&\approx U(\theta_0)-H(\theta_0)(\widehat\theta-\theta_0)
                   -K\theta_0-K(\widehat\theta-\theta_0),\\
 (H+K)(\widehat\theta-\theta_0)&\approx U(\theta_0)-K\theta_0 .
 \label{cov-penalty-linearization}
\end{align}
Replacing the random sensitivity by its suitable population counterpart
$A_0$ identifies both pieces:
\[
 \E(\widehat\theta-\theta_0)
       \approx-(A_0+K)^{-1}K\theta_0,\qquad
 \Var(\widehat\theta)\approx
       (A_0+K)^{-1}J_0(A_0+K)^{-\top}.
\]
The first expression describes local smoothing bias.  When the penalty
does not become negligible, a population penalized target can differ
substantially from $\theta_0$, and linearization should instead be
performed at that target.  The displayed local bias is not a globally
exact formula for a nonlinear count regression.

A plug-in penalized sandwich uses $H+K$ as bread and an appropriate
estimate of the variability of the \emph{data} score as meat.  The
deterministic penalty is not an independent random contribution and
must not be added to every study's score.  Under strong shrinkage,
data-score outer products may also require attention to centering and
the definition of the penalized population target; an uncentered meat
is not universally consistent when individual score means are nonzero.
Under a correctly specified likelihood, the local expected-curvature
version is $(I+K)^{-1}I(I+K)^{-1}$.

By contrast, a Gaussian prior with precision $K$ gives negative
log-posterior curvature $H+K$.  A Laplace posterior approximation uses
$(H+K)^{-1}$.  It is exactly the posterior covariance only for an
actually quadratic posterior with appropriate support.  A singular
roughness precision corresponds to unpenalized directions and need
not define a proper prior on its own.

The expected-curvature versions differ by the exact identity
\begin{align}
 (I+K)^{-1}-(I+K)^{-1}I(I+K)^{-1}
 &=(I+K)^{-1}\{(I+K)-I\}(I+K)^{-1}\\
 &=(I+K)^{-1}K(I+K)^{-1}\succeq0,
 \label{cov-penalty-covariance-gap}
\end{align}
assuming $I+K\succ0$ and $K\succeq0$.
For a one-dimensional quadratic example with information $a>0$ and
penalty $k\geq0$, inverse curvature is $1/(a+k)$, whereas the
sampling variance is $a/(a+k)^2$ and the local bias is
$-k\theta_0/(a+k)$.  Neither variance includes the squared bias
needed for mean squared error about $\theta_0$.

\subsection{Selected smoothing parameters, null spaces, and numerical ridge}
\label{cov-penalty-selection}

All previous penalty derivatives hold $\lambda$ fixed.  If $\lambda$
is selected from the same data, uncertainty in $\widehat\lambda$
and its dependence on the regression estimate are omitted by a
fixed-$\lambda$ calculation.  For a smooth solution branch,
differentiating $U(\widehat\theta)-K(\lambda)\widehat\theta=0$
gives
\[
 (H+K)\frac{\partial\widehat\theta}{\partial\lambda_a}
       =-\frac{\partial K}{\partial\lambda_a}\widehat\theta .
\]
A joint delta calculation must include the variability of
$\widehat\lambda$ and the cross covariances with the fixed-smoothing
fit, not merely append an independent variance term.  A bootstrap
that repeats the selection rule is another possible route under its
own assumptions.

When $H\succeq0$ and $K\succeq0$,
\[
 v^\top(H+K)v=0
 \quad\Longleftrightarrow\quad v\in\ker(H)\cap\ker(K).
\]
The implication uses the nonnegativity of both terms and the spectral
characterization of the null space of a PSD matrix.  Thus
$H+K\succ0$ precisely when their null spaces have no common nonzero
direction.  A roughness penalty that leaves constant maps unpenalized
cannot resolve a likelihood confounding direction that is also constant.
For an indefinite $H$, this PSD null-space argument does not apply.

For example, if $Bc_B=\mathbf1_N$ and $z_i^\top c_\gamma=1$
for every study, then shifting every group map by $tc_B$ and
shifting $\gamma$ by $-tc_\gamma$ leaves the predictor unchanged.
Likelihood identification requires a constraint or removal of the
redundant direction.  Penalizing that direction can select a unique
answer, but it introduces additional assumptions rather than new data
information.

Adding a numerical ridge $\varepsilon I_{n_\theta}$ solely during
covariance calculation changes the inverse being reported without
changing the estimating equation that produced the fit.  It is not
automatically a fitted penalty or a sampling covariance for that fit.
The repository's \texttt{ridge} option is explicitly numerical
regularization.  Its magnitude and inferential effect should be
reported rather than described as magical identification.

\section{Joint nuisance dispersion and reparameterization}
\label{cov-dispersion}

\subsection{Regression uncertainty after estimating dispersion}
\label{cov-nuisance-schur}

Now let $\alpha\in(0,\infty)^L$ collect dispersion parameters; $L=1$ is
allowed.  Use a distinct joint information matrix
\[
 \mathcal H=
 \begin{pmatrix}
 H_{\theta\theta}&H_{\theta\alpha}\\
 H_{\alpha\theta}&H_{\alpha\alpha}
 \end{pmatrix}.
\]
Here $H_{\theta\theta}$ is the regression curvature with dispersion held
fixed at its evaluation value, not the inverse-covariance block of a
joint fit.  If the nuisance block is nonsingular, eliminate the
nuisance variables to define
\begin{equation}
 H_{\mathrm{eff}}=
 H_{\theta\theta}
   -H_{\theta\alpha}H_{\alpha\alpha}^{-1}H_{\alpha\theta},
 \qquad
 (\mathcal H^{-1})_{\theta\theta}=H_{\mathrm{eff}}^{-1}.
 \label{cov-nuisance-adjustment}
\end{equation}
This is the same Schur identity with the roles of the block variables
reordered.  If $\mathcal H\succ0$, the subtraction is PSD and
$H_{\mathrm{eff}}^{-1}\succeq H_{\theta\theta}^{-1}$.
For one nuisance parameter with cross column $h$ and scalar
curvature $a>0$, the effective information is
$H_{\theta\theta}-hh^\top/a$.  A small nuisance dimension does not
guarantee a small adjustment.

Expected information may have zero regression--dispersion cross blocks
for some parametrizations and models.  That must be established for
the particular likelihood; it cannot be inferred from alternating
optimization, and it does not force the observed cross block to be
zero.  For robust inference, stack the regression and dispersion
estimating equations, form their joint bread and joint score meat,
and then project the full sandwich onto regression coordinates.
Replacing a regression bread by a Schur complement alone while
ignoring nuisance-score variation is not a general robust adjustment.

The regression information and covariance interfaces discussed here in
\texttt{cbmr\_improved} hold the fitted dispersion fixed.  They do
not supply the joint nuisance adjustment in
Equation~\eqref{cov-nuisance-adjustment}.  A likelihood fit that updates
dispersion is not, by that fact alone, accompanied by joint-dispersion
standard errors.

\subsection{Log dispersion and the extra Hessian term}
\label{cov-log-dispersion}

For positive dispersions write $\delta_a=\log\alpha_a$, so
$\alpha_a=e^{\delta_a}$, and put $A_\alpha=\diag(\alpha)$.
Let $U_\alpha=\nabla_\alpha\ell$ at the evaluation point.  The
first derivative transforms as
$U_\delta=A_\alpha U_\alpha$.  Differentiating its $a$th component
with respect to $\delta_b$ gives
\[
 \frac{\partial (U_\delta)_a}{\partial\delta_b}
 =\mathbf1\{a=b\}\alpha_a(U_\alpha)_a
   +\alpha_a\alpha_b
      \frac{\partial^2\ell}{\partial\alpha_a\partial\alpha_b}.
\]
Changing sign to information therefore yields
\begin{equation}
 H_{\delta\delta}
   =A_\alpha H_{\alpha\alpha}A_\alpha
            -\diag\bigl(\alpha_a(U_\alpha)_a:a=1,\ldots,L\bigr),
 \qquad H_{\theta\delta}=H_{\theta\alpha}A_\alpha .
 \label{cov-log-hessian}
\end{equation}
For a scalar this reads
$H_{\delta\delta}=\alpha^2H_{\alpha\alpha}-\alpha U_\alpha$.
The minus sign of the extra score term follows from the convention
$H=-\nabla^2\ell$.

At an interior joint stationary point $U_\alpha=0$, the extra term
vanishes and, with
$J_{\mathrm{par}}=\blkdiag(I_{n_\theta},A_\alpha)$,
the joint observed information transforms by
$\mathcal H_{(\theta,\delta)}
=J_{\mathrm{par}}^\top\mathcal H_{(\theta,\alpha)}J_{\mathrm{par}}$.
Expected information obeys the corresponding rule under zero expected
score and regularity.  Away from stationarity, simply applying this
congruence to an observed Hessian drops a real chain-rule term.
At a penalized optimum the likelihood score alone may be nonzero:
transform the full objective and its gradient if that is the curvature
being used.  Finally, $\alpha=0$ is a boundary not represented by
finite $\delta$; interior Wald arguments do not automatically cover it.

\section{Voxel contrasts, nonlinear maps, and inferential calibration}
\label{cov-inference}

\subsection{SV coefficient ordering and tensor-product contrasts}
\label{cov-sv-contrasts}

Let $C_Z\in\R^{k\times R}$ specify contrasts among SV maps.
At voxel $j$, stacking the $R$ log effects gives
\[
 \begin{pmatrix}B_j^\top\beta_1\\ \vdots\\B_j^\top\beta_R\end{pmatrix}
 =(I_R\otimes B_j^\top)\beta .
\]
Multiplication by $C_Z$ and the Kronecker mixed-product identity give
\begin{equation}
 \widehat\delta_j
   =(C_Z\otimes B_j^\top)\widehat\beta-\delta_{0j},\qquad
 V_{\delta j}
   =(C_Z\otimes B_j^\top)V_{\beta\beta}
                    (C_Z^\top\otimes B_j).
 \label{cov-sv-contrast-formula}
\end{equation}
The contrast matrix is $k\times RP$, not $RP\times k$.
If global coefficients are present, append $k\times Q$ zeros to
obtain a contrast on $\theta$.  For example, a row with entries
$(1,-1,0,\ldots,0)$ produces the coefficient row
$[\,B_j^\top,-B_j^\top,0,\ldots,0\,]$.
This explicitly checks that the stacking convention agrees with the
tensor-product order.  Contrast row sums need not equal one; differences
usually sum to zero.

\subsection{GC group differences and global tests}
\label{cov-gc-contrasts}

With $C_G\in\R^{k\times G}$, use
$\mathcal T_{Gj}=[\,C_G\otimes B_j^\top\ \ 0_{k\times Q}\,]$.
Then $\widehat\delta_j=\mathcal T_{Gj}\widehat\theta-\delta_{0j}$
and $V_{\delta j}=\mathcal T_{Gj}V\mathcal T_{Gj}^\top$,
where $V$ must be the covariance choice appropriate to the estimator.
For one difference $\xi_g-\xi_h$, expansion gives
\begin{align}
 \Var\{B_j^\top(\widehat\xi_g-\widehat\xi_h)\}
 &=B_j^\top(V_{gg}+V_{hh}-V_{gh}-V_{hg})B_j\\
 &=B_j^\top(D_g^{-1}+D_h^{-1})B_j\\
 &\quad+B_j^\top(T_g-T_h)S^{-1}(T_g-T_h)^\top B_j,
 \label{cov-gc-difference}
\end{align}
where the last two lines assume the observed inverse-curvature choice
and $g\ne h$.  This expression retains shared-global uncertainty and
shows when it cancels from a difference.  Unequal study numbers
per group affect the information but do not invalidate the identity.

A global contrast $C_\gamma\gamma$ has local model-based covariance
$C_\gamma S^{-1}C_\gamma^\top$, not
$C_\gamma\mathcal E^{-1}C_\gamma^\top$.  A comparison of fitted
intensities at different global covariate values must include the
appropriate nonzero global contrast as well; it is not necessarily a
pure group-map contrast.

\subsection{Exponentiation and normalized intensities by the delta method}
\label{cov-delta}

For one map let $\eta=B\xi$ and $\mu=\exp(\eta)$ componentwise.
Since $\partial\mu_j/\partial\xi^\top=\mu_jB_j^\top$,
the exact Jacobian is $J_\mu=\diag(\mu)B$.  The first-order expansion
\[
 \mu(\widehat\xi)-\mu(\xi_0)
       \approx J_\mu(\xi_0)(\widehat\xi-\xi_0)
\]
leads to
$\widehat{\Cov}\{\mu(\widehat\xi)\}
\approx J_\mu(\widehat\xi)V_{\xi\xi}J_\mu(\widehat\xi)^\top$.
The Jacobian is exact, whereas the covariance approximation is not.
For a scalar log contrast $\delta$, the ratio $r=e^\delta$ has
approximate variance $r^2\Var(\widehat\delta)$.  Exponentiating a
log-scale interval often respects positivity better than a symmetric
interval on the ratio scale.

For normalized spatial intensity define
$p_j=\mu_j/m_\mu$, where $m_\mu=\sum_l\mu_l$.
The quotient rule gives the entire normalization derivative:
\begin{align}
 dp_j
 &=\frac{m_\mu\,d\mu_j-\mu_j\sum_l d\mu_l}{m_\mu^2}\\
 &=p_j\left(d\eta_j-\sum_lp_l\,d\eta_l\right),\\
 J_p&=\{\diag(p)-pp^\top\}B .
 \label{cov-normalized-jacobian}
\end{align}
Thus $\Cov\{p(\widehat\xi)\}\approx J_pV_{\xi\xi}J_p^\top$.
Because $\mathbf1_N^\top J_p=0$, this covariance is singular in the
sum-to-one direction and has rank at most $\min(P,N-1)$.
A common multiplicative global intensity factor cancels from this
within-map normalization exactly.  This does not mean global
estimation has no indirect effect on the covariance of the fitted
spatial coefficients.  Covariances between normalized group maps
still require $J_{p,g}V_{gh}J_{p,h}^\top$.

\subsection{Wald tests: algebra versus the limiting null law}
\label{cov-wald}

For a $k$-dimensional contrast with nonsingular estimated covariance,
the Wald statistic is computed by solving $V_{\delta j}a=\widehat\delta_j$
and taking
\[
 W_j=\widehat\delta_j^\top a
     =\widehat\delta_j^\top V_{\delta j}^{-1}\widehat\delta_j .
\]
If the null is true, the estimator is asymptotically normal in the
tested directions, the covariance estimate is consistent, and the
$k$ constraints are independent, then $W_j$ converges in distribution
to $\chi_k^2$.  For one contrast the signed statistic is
$Z_j=\widehat\delta_j/\sqrt{V_{\delta j}}$, with $W_j=Z_j^2$.
These distributional statements are not exact consequences of Schur
inversion.

Redundant constraints should be reduced to an independent set.
A generalized inverse can describe a known lower-dimensional
estimable test under additional rank conditions, but cannot rescue
unidentified contrasts, too few independent clusters, or an unstable
covariance.  Smoothing bias, estimated nuisance parameters, boundary
fits, and data-dependent model selection all require attention.
Pointwise voxel $p$-values are not automatically multiplicity-adjusted
or simultaneous spatial inference.  Faster exact linear algebra does
not establish a universal minimum number of foci or studies for Wald
calibration.

\subsection{Bootstrap units, null fitting, and repeated fitting decisions}
\label{cov-bootstrap}

A bootstrap must reproduce the estimator and the relevant sampling
mechanism.  For a nonparametric study bootstrap, resample independent
studies (or larger independent clusters), retaining each selected
unit's entire voxel pattern and covariates.  Group-stratified resampling
may be appropriate when the design conditions on group counts.
Resampling individual voxels instead changes the within-study dependence
model.  A fixed-design parametric bootstrap instead keeps covariates
and group membership fixed and simulates new responses.

For a fitted-null parametric test, the operational sequence is:
fit the stated null with its constraints; simulate under that null
and its dependence model; refit each simulated dataset by the original
procedure; and recompute the same statistic.  If the procedure estimates
dispersion or selects smoothing parameters, repeat those steps unless
the explicit target is conditional on their fixed values.  A shared
study latent effect must be shared across that study's simulated voxels.
Independent-cell NB and aggregate-NB sampling schemes are not
interchangeable with a study-clustered NB scheme.

With $B_{\mathrm{boot}}$ successful replicates, a common nonzero
Monte Carlo tail estimate is
\[
 \widehat p=
 \frac{1+\sum_{b=1}^{B_{\mathrm{boot}}}
                     \mathbf1\{W_b^*\geq W_{\mathrm{obs}}\}}
      {B_{\mathrm{boot}}+1}.
\]
The statistic and tail must match the alternative, and failed refits
must be diagnosed rather than silently discarded as if random.
The $+1$ convention does not make a fitted-null parametric bootstrap
an exact finite-sample test.  Bootstrap validity and any tail
extrapolation remain separate statistical questions.

\section{Conditioning, Kronecker structure, and computational accounting}
\label{cov-computation}

\subsection{Deriving the spectral condition number}
\label{cov-condition}

For a real symmetric nonsingular $H$, take the orthogonal
eigendecomposition $H=O\Lambda O^\top$, with
$\Lambda=\diag(\lambda_1,\ldots,\lambda_{n_\theta})$.
Then
\[
 H^\top H=O\Lambda^2O^\top .
\]
The singular values are the square roots of the eigenvalues of this
product, namely $|\lambda_a|$.  Therefore
\begin{equation}
 \|H\|_2=\max_a|\lambda_a|,\qquad
 \|H^{-1}\|_2=\frac1{\min_a|\lambda_a|},\qquad
 \kappa_2(H)=\frac{\max_a|\lambda_a|}{\min_a|\lambda_a|}.
 \label{cov-spectral-condition}
\end{equation}
For SPD $H$, absolute values can be dropped.  For an indefinite matrix,
using the largest signed eigenvalue divided by the smallest signed
eigenvalue is wrong.  A zero eigenvalue makes the matrix singular and
the ordinary condition number infinite.  The example
$\diag(1,-1)$ has $\kappa_2=1$ but cannot be a covariance precision
for a nondegenerate Gaussian distribution.

For a block diagonal matrix,
$\det(tI-D)=\prod_g\det(tI-D_g)$, so its spectrum is the union of
block spectra and
\[
 \kappa_2(D)=
 \frac{\max_g\max_a|\lambda_a(D_g)|}
      {\min_g\min_a|\lambda_a(D_g)|}.
\]
This is not generally $\max_g\kappa_2(D_g)$.  For example,
$D_1=I_P$ and $D_2=10^6I_P$ each have condition number one,
whereas $\blkdiag(D_1,D_2)$ has condition number $10^6$.
The scale across blocks matters even when each block is internally
well conditioned.

For a bordered system, Equation~\eqref{cov-congruence} preserves
inertia, not eigenvalues.  In the scalar-block example
$D=1$, $F=t$, $\mathcal E=1+t^2$, the Schur complement is $S=1$.
Both factors have condition number one, but the full matrix has
determinant one and trace $t^2+2$.  Its eigenvalues are
\[
 \lambda_\pm=\frac{t^2+2\ \pm\sqrt{(t^2+2)^2-4}}2,
\]
so $\kappa_2(H)=\lambda_+/\lambda_-$ grows without bound as
$|t|$ increases.  A full symmetric eigenproblem is required for
the exact bordered spectral condition number.

Conditioning depends on coefficient units.  Under a nonsingular linear
reparameterization $\theta=L\phi$, the information becomes
$H_\phi=L^\top H_\theta L$ and can have a different condition number,
while correctly transformed contrasts describe the same estimand.
For fixed $H$, perturbing a right-hand side gives the bound
$\|\Delta x\|_2/\|x\|_2
\leq\kappa_2(H)\|\Delta b\|_2/\|b\|_2$.
This follows from $\Delta x=H^{-1}\Delta b$ and
$\|b\|_2\leq\|H\|_2\|x\|_2$; it explains numerical sensitivity,
not statistical validity.  A one-norm estimator such as LAPACK's
\texttt{pocon} estimates a different norm condition number and is
not an exact computation of $\kappa_2$.

\subsection{When a Kronecker inverse is justified}
\label{cov-kronecker}

For the SV independent-cell information without a global border,
Equation~\eqref{cov-sv-block} can be written
\[
 H=\sum_{i,j}w_{ij}(Z_iZ_i^\top)\otimes(B_jB_j^\top).
\]
If the weights really separate as $w_{ij}=a_i b_j$, distributivity
of the tensor product gives the exact factorization
\begin{equation}
 H=A_Z\otimes A_B,\qquad
 A_Z=\sum_i a_iZ_iZ_i^\top,\quad
 A_B=\sum_j b_jB_jB_j^\top .
 \label{cov-kron-factorization}
\end{equation}
This is a sufficient separability condition, not a consequence of
having written an unweighted design as a Kronecker product.  Generic
SV means vary with both study and voxel through several interacting
maps.  Independent-NB denominators and observed curvature weights
can destroy a mean factorization as well.  A sum of Kronecker products
is not generally a single Kronecker product.  A global border or a
penalty also requires an explicit new structural check.

If the two square factors are nonsingular, the mixed-product rule
verifies the inverse:
\[
 (A_Z\otimes A_B)(A_Z^{-1}\otimes A_B^{-1})
 =(A_ZA_Z^{-1})\otimes(A_BA_B^{-1})=I_R\otimes I_P.
\]
The reverse product is the identity too.  If either factor is singular,
a nonzero null vector in that factor produces a null vector of the
product, so an ordinary inverse is unavailable.  Singular values
of a Kronecker product are pairwise products of singular values of
its factors; consequently, for nonsingular factors,
\[
 \kappa_2(A_Z\otimes A_B)=\kappa_2(A_Z)\kappa_2(A_B).
\]
For symmetric factors this also follows from tensor products of their
orthonormal eigenvectors and products of absolute eigenvalues.

To solve without forming the product, reshape a right-hand side as
$b=\operatorname{vec}(Y)$ with $Y\in\R^{P\times R}$, stacking
columns in the same map-major order as $\beta$.  The identity
\[
 (A_Z\otimes A_B)\operatorname{vec}(X)
       =\operatorname{vec}(A_BXA_Z^\top)
\]
turns the equation into $A_BXA_Z^\top=Y$.  First solve
$A_BW=Y$, then solve $A_ZX^\top=W^\top$.
Using an approximate separable matrix as a preconditioner can be useful,
but its inverse is not the exact covariance of the original information.

\subsection{Operation counts, storage, and implementation scope}
\label{cov-costs}

The following counts concern linear algebra \emph{after} information
assembly, use ordinary dense block arithmetic, and omit constant factors.
Let block $g$ have size $d_g$, with $\sum_gd_g=d$; use component
indices in place of group indices for general SV designs.

\begin{center}
\begin{tabularx}{\linewidth}{@{}lXX@{}}
\toprule
Representation & Initial factorization work & Factor storage\\
\midrule
Full dense, $n_\theta=d+Q$
 & $O(n_\theta^3)$
 & $O(n_\theta^2)$\\
Unbordered blocks
 & $O(\sum_gd_g^3)$
 & $O(\sum_gd_g^2)$\\
Bordered blocks
 & $O(\sum_gd_g^3+Q\sum_gd_g^2+dQ^2+Q^3)$
 & $O(\sum_gd_g^2+dQ+Q^2)$\\
Exact $A_Z\otimes A_B$
 & $O(R^3+P^3)$
 & $O(R^2+P^2)$\\
\bottomrule
\end{tabularx}
\end{center}

For the bordered row, the terms arise in order: factor every $D_g$;
solve $D_gT_g=F_g$ for $Q$ right-hand sides; form
$S=\mathcal E-F^\top T$; and factor $S$.
Afterwards, solving $k$ general right-hand sides costs
\[
 O\left(k\sum_gd_g^2+kdQ+kQ^2\right).
\]
The block triangular solves account for the first term,
$F^\top y$ and $Tx_\gamma$ for the second, and the Schur solves
for the third.  A full dense solve after factorization costs
$O(kn_\theta^2)$.  For the Kronecker system, the two factor solves
cost $O(P^2R+PR^2)$ per right-hand side.  Holding its reshaped
right-hand side requires $O(PR)$ additional storage.

Selected covariance output has its own costs.  Storing $X$ in
Equation~\eqref{cov-selected-solve} requires $O(n_\theta k)$,
and a dense multiplication $CX$ costs $O(n_\theta k^2)$.
If the full spatial correction is demanded, solving
$SW=T^\top$ costs $O(Q^2d)$ and forming $TW$ costs $O(d^2Q)$;
storing the resulting full covariance requires $O(n_\theta^2)$.
Computing all diagonal entries avoids that output allocation, though
generic block-inverse diagonals may still need
$O(\sum_gd_g^3)$ work via solves against block identity matrices.

For projected study sandwiches, forming $\mathsf U X$ costs
$O(Mn_\theta k)$ with dense score rows and then its Gram matrix
costs $O(Mk^2)$.  If scores arrive study by study, one can accumulate
$v_iv_i^\top$ immediately and retain only $O(k^2)$ output storage,
in addition to the factors and the contrast solves.  Model-specific
sparse projections can reduce score-assembly work.  These are order
calculations, not empirical performance claims; likelihood assembly,
input storage, sparsity, and numerical libraries affect actual timings.

The SPD path in \texttt{cbmr\_improved/structured.py} caches block
Cholesky factors, computes $T$ by solves, and factorizes the Schur
complement.  Its contrast, diagonal, and score-Gram operations implement
the algebra above without a full spatial inverse on that path.
If block or Schur Cholesky fails, it logs a fallback to a dense symmetric
solve; an invertible indefinite matrix can then be solved algebraically,
while a singular full system raises an error.  The structured savings
need not survive this fallback.  Its full covariance compatibility
method deliberately solves against the full identity.

The same implementation computes the spectral condition number from
block eigenvalues when there is no global border, and from a full
dense symmetric eigenproblem when there is one.  That diagnostic itself
can require $O(n_\theta^2)$ memory and $O(n_\theta^3)$ work even if
the intended contrast solve is structured.  Exact algebraic identities
do not remove floating-point error, establish statistical calibration,
or constitute formal certification of this explanatory fragment.

\subsection{Manual audit prompts}
\label{cov-manual-audit}

\begin{enumerate}
 \item Which regression coefficients and nuisance parameters were actually
 estimated?  Does the stated covariance condition on any of them?
 \item Is the matrix observed information, expected information,
 penalized curvature, or an estimating-equation sensitivity?  Are its
 sign and summed-versus-averaged normalization explicit?
 \item Can every claimed structural zero be traced to factor support or a
 specified penalty?  Are entire connected components, not isolated
 missing edges, used to choose blocks?
 \item Do the products $DT=F$, $S=\mathcal E-F^\top T$, and
 $HX=C^\top$ have small residuals at the numerical fit?  Are the
 matrices whose positive definiteness is required actually SPD?
 \item Are group cross covariances retained, and does a selected contrast
 use the same coefficient ordering as $C_Z\otimes B_j^\top$ or
 $C_G\otimes B_j^\top$?
 \item For a sandwich, are the score rows derivatives of the chosen
 likelihood or estimating equation, grouped by genuinely independent
 units?  Is their rank sufficient for the proposed test?
 \item For a penalty or numerical ridge, what changes in the estimator,
 bias, covariance interpretation, and identification assumptions?
 Was smoothing selection uncertainty intentionally included or omitted?
 \item Does a proposed Kronecker factorization hold for the actual
 weights at this fit, rather than only for the unweighted design?
 \item Which statements are exact matrix identities, which are
 first-order approximations, and which rely on a bootstrap or empirical
 calibration?  Have finite-sample validity and multiplicity been assessed
 separately from computational speed?
\end{enumerate}
